\section{Preliminaries}
Consider a generic multi-matrix model in $n$ Hermitian matrix degrees of freedom $\mathbf{X}=\{X_i\}_{i=0}^{n-1}$ of the form
%
\begin{equation}
    \label{loops1}
    Z = \int \prod_{i=0}^{n-1} \rmd X_i \, e^{-N\Tr V(\mathbf{X})},
\end{equation}
%
where $V(\mathbf{X})$ is a potential function which is a polynomial of degree $d+1$ in the $n$ matrices of $\mathbf{X}$, and depends on couplings $\{t_i\}_{i=2}^{\infty}$. We are interested in computing the planar resolvent of a given matrix, say $X = X_0$,
\beq\label{res}
W(z) = \frac{1}{N}\left\langle \Tr \frac{1}{z-X}\right\rangle
\eeq
\noindent where here and in what follows $\langle\cdot\rangle$ denotes the average with respect to the measure \eref{loops1} and $z$ denotes the fugacity of a boundary link.

Define the following resolvent functions with boundary insertions
%
\beq\label{loopRes}
    W_{w(\mathbf{X})}(z) = \frac{1}{N} \left\langle \Tr  w(\mathbf{X}) \frac{1}{z-X} \right\rangle,
\eeq
%
where $w(\mathbf{X})$ is a word in the free algebra generated by $\mathbf{X}$, but with the constraint that $X$ is not the first or last letter of the word. On the left hand side we will proceed to label the resolvent functions by the indices labelling the matrices that comprise the word $w(\mathbf{X})$ for the sake of brevity. Hence, one may write
%
\beq\label{resExample}
    W_{(1,0,2)}(z) = \frac{1}{N} \left\langle \Tr X_1 X_0 X_2\frac{1}{z-X} \right\rangle.
\eeq
%
Furthermore, we define the following trace correlation functions analogously,
%
\begin{equation}
    \label{trace}
    p_{w(\mathbf{X})}= \frac{1}{N} \left\langle \Tr w(\mathbf{X})\right\rangle.
\end{equation}
%
Define the level of a resolvent function of the form \eref{loopRes} as the length of the associated word, $L(w(\mathbf{X}))$, such that, for example, \eref{resExample} with $L(X_1 X_0 X_2)=3$ denotes a level-3 resolvent function. Similarly, we index trace correlation functions of a given level by the set of words of appropriate length. 

The number of inequivalent trace correlation functions of a given level is generally much smaller than the size of the corresponding set of words. Due to cyclicity of the trace there is a $D_k$ symmetry for the level-$k$ correlators, where $D_k$ is the dihedral group of order $k$. There may also be a global symmetry, $G$, of the matrix model potential, given by
%
\beq
    \label{symmetry}
    G = \{ g\in \mathrm{Aut}(\mathbf{X}) : V(g(\mathbf{X}))=V(\mathbf{X})\}.
\eeq
%
An example of this would be permutation symmetry amongst the matrix degrees of freedom. Therefore, defining the set of irreducible trace correlation functions of level $k$ as $\mathcal{P}_k$, we may write
%
\begin{equation}
    \label{irreducibleTrace}
    \mathcal{P}_k \cong \mathbf{X}^k / (D_k \times G).
\end{equation}
%
The number of inequivalent correlators is then given by the cardinality of this set.

Similarly we define the set of irreducible resolvent functions of level $k$ as $\mathcal{W}_k$. The number of independent resolvent functions is also constrained, although not as severely as the correlation functions are. Define the operator $\Delta$, acting on resolvent functions as follows:
%
\begin{eqnarray}
    \label{delta}
    \Delta^k W_{w(\mathbf{X})}(z) & : = & \frac{1}{N} \left\langle \Tr  X^k w(\mathbf{X}) \frac{1}{z-X}\right\rangle\\
    & = & z^k W_{w(\mathbf{X})}(z) - \sum_{i=1}^k z^{k-i} p_{X^{i-1}w(\mathbf{X})}.
\end{eqnarray}
%
We may omit words for which any sub-word consisting only of the matrix $X$ adjoins the loop operator $(z-X)^{-1}$, since these may be expressed in terms of the $\Delta$ operator to a power given by the length of the adjoining sub-word in $X$, acting on a lower-level resolvent function. Moreover, the presence of the loop operator means that the dihedral symmetry is removed. Finally, having fixed $X$ as the representative matrix of the loop operator, the defining word only respects the subgroup $G'\subseteq G$, consisting of elements $g\in\mathrm{Aut}(\tilde{\mathbf{X}})$, where $\tilde{\mathbf{X}}=\mathbf{X} \setminus X$. Therefore we may write
%
\begin{equation}
    \label{irreducibleResolvent}
    \mathcal{W}_k \cong (\tilde{\mathbf{X}} \times \mathbf{X}^{k-2} \times \tilde{\mathbf{X}})/G'.
\end{equation}

For example, consider a three-matrix model with matrix degrees of freedom $X_0,X_1,X_2$, with a $\mathbb{Z}_2$ symmetry under exchange of $X_0$ and $X_1$. Then for the resolvent functions with respect to $X_1$, one would find $ W_{(2,1,0,1,2)} \neq W_{(2,0,1,0,2)}$, and therefore the representative words comprise two distinct elements of $\mathcal{W}_5$,
whereas the same words are identified in $\mathcal{P}_5$ as one has $p_{2,1,0,1,2} = p_{2,0,1,0,2}$.

Loop equations are Schwinger-Dyson equations of the matrix model, expressing the invariance of the matrix integral under infinitesimal reparameterisations $X_i \rightarrow X_i + \varepsilon \delta X_i$. There is some freedom as to which variations one can consider. For example, in the approach of \cite{Eynard:2004mh}, specific to the two-matrix model, one computes reparameterisations using products of the resolvent operators for the different matrices. The approach we adopt here is to consider variations of the form
%
\begin{equation}
    \label{reparam}
    \delta X_i = w_1(\mathbf{X}) \frac{1}{z-X} w_2(\mathbf{X}).
\end{equation}
%
One then evaluates the loop equation by equating the first order contribution from the Jacobian with the contribution from the potential. For a variation of $X=X_0$, differentiating the resolvent gives the Jacobian contribution prescribed by the `split rule' \cite{Eynard:2004mh},
\beq
\label{Jsplit}
J_{\mathrm{res}}(X)=\Tr \left( w_1(\mathbf{X}) \frac{1}{z-X}\right)\Tr\left(\frac{1}{z-X} w_2(\mathbf{X}) \right).
\eeq
There are additional Jacobian terms from differentiating any occurrences of the varied matrix in $w_1$ and $w_2$. For variations of $X_i$ with $i\neq0$, only these inserted-word terms contribute, since the resolvent is independent of $X_i$.
The resulting constraints may be expressed in terms of the resolvent functions \eref{loopRes} and trace correlation functions \eref{trace}. 

We can classify loop equations by the level of their highest-level resolvent functions. If $V'(\mathbf{X})$ is degree $d$ in $\tilde{\mathbf{X}}$\footnote{The degree will generally be different for variations corresponding to different $X_i\in\mathbf{X}$.}, then level-$k$ loop equations are given by \eref{reparam} with $L(w_2)=k-d-L(w_1)$. We can therefore define a corresponding set of level $k$ variations,
%
\begin{equation}
    \label{var}
    \mathcal{V}_k = \left\{ (w_1,w_2)\in\mathbf{V}_m\times\overline{\mathbf{V}}_{k-d-m} : m=k-d,\cdots, \left\lceil\frac{k-d}{2}\right\rceil \right\},
\end{equation}
%
where $\mathbf{V}_m=\mathbf{X}^{m-1}\times \tilde{\mathbf{X}}$ and $\overline{\mathbf{V}}_m = \tilde{\mathbf{X}}\times\mathbf{X}^{m-1}$. That is, the words $w_1$ and $w_2$ are constrained such that the letters in contact with the loop operator are in $\tilde{\mathbf{X}}$. This is so that the resulting loop equations are distinct from lower-level loop equations.

The key idea of the loop equation approach we adopt is to compute the Schwinger-Dyson constraints for each matrix degree of freedom for a sufficiently large set of variations. We aim to compute loop equations up to a level $n_{\mathrm{max}}$ with enough independent constraints to eliminate the auxiliary resolvent functions and retain an algebraic equation for the planar (level-0) resolvent. Comparing the number of equations with the number of auxiliary resolvent functions to be eliminated provides a useful guide to the required level. This count alone does not establish closure: dependencies among the equations, relations among the correlator moduli, and special values of the couplings can change the effective rank. Closure must therefore be checked by the elimination itself. This may be regarded as a discrete version of loop equation solution to the two-matrix model \cite{Eynard:2004mh}, since one may expand the product of resolvent operators as an infinite sum of words in matrices of arbitrary length multiplying a single resolvent operator. It is hoped that this refined set of equations contains more information than if one only considers a small set of reparameterisations by products of loop operators, allowing us to compute the planar resolvent.
